\section[0. AI Project]{0. AI Project}

\subsection{Intruduction}

\begin{itemize}
\item \textbf{인공지능}\textit{\textsuperscript{Artifitial Intelligence}}: 컴퓨터가 사람처럼 생각하고 판단하게 만드는 기술
\item \textbf{머신러닝}\textit{\textsuperscript{Machine Learning}}: 인공지능의 한 분야. 인간의 학습 능력과 같은 기능을 컴퓨터에 부여하기 위한 기술
\item \textbf{딥러닝}\textit{\textsuperscript{Deep Learning}}: 인공신경망 기반으로 한 머신러닝의 방법 중 하나. 빅데이터 기반으로 스스로 학습하여 판단하는 기술
\end{itemize}

파이썬으로 인공지능을 코딩하는 것은 관련한 다양한 라이브러리와 같은 여러 요인들로 인해 간편하다는 장점이 있다.

\subsection{Machine Learning Type}

\begin{itemize}
\item \textbf{지도 학습}\textit{\textsuperscript{Supervised Learning}}: 문제와 정답을 알려주고 공부. (키워드: 예측, 분류)
\item \textbf{비지도 학습}\textit{\textsuperscript{Unsupervised Learning}}: 답을 가르쳐주지 않고 공부. (키워드: 연관 규칙, 군집)
\item \textbf{강화 학습}\textit{\textsuperscript{Reinforcement Learning}}: 보상을 통해 상을 최대화, 벌을 최소화하는 방향으로 학습. (키워드: 보상)
\end{itemize}

\subsection{Examples}

\begin{itemize}
\item Computer Vision: 컴퓨터가 이미지나 비디오 해석
    \begin{itemize}
    \item \textbf{Object Detection}: 이미지에서 의미 있는 객체의 종류와 위치(bounding box) 찾음
    \item \textbf{Object Tracking}: 비디오에서 움직이는 객체를 식별하고 추적
    \item \textbf{Face Recognition}: 이미지, 비디오에서 얼굴을 인식하고 식별
        \begin{itemize}
        \item 보로노이 알고리즘: 얼굴에 점을 찍어서 얼굴 인식함.
        \end{itemize}
    \item \textbf{Image Generation}: 텍스트나 다른 이미지 기반으로 새로운 이미지 생성
    \end{itemize}
\item NLP: 컴퓨터가 인간의 언어를 이해하고 생성함
    \begin{itemize}
    \item \textbf{Machine Translation}: 번역
    \item \textbf{Chatbot and Conversational System}: 챗봇 및 대화. chatbot은 데이터가 부족하여 영화 대본을 기반으로 만듦.
    \item \textbf{Sentiment Analysis}: 감정 분석
    \item \textbf{Text Summarization}: 텍스트 요약
    \end{itemize}
\end{itemize}

AI 활용 예시: 헬로미러, 얼굴 인식 인공지능 스피커, 자율주행 및 장애물 회피(POSLA, RC Car, 자동 쇼핑카트), 군집 주행 드론, 주가 예측, 시험 부정 방지, 메타버스 토론
